\documentclass[11pt,a4paper]{article}
\usepackage{cmap}
\usepackage[T2A]{fontenc}
\usepackage[utf8]{inputenc}
\usepackage[russian]{babel}
\usepackage{amsmath,amssymb,amsthm}
\usepackage{geometry}
\geometry{left=20mm,right=20mm,top=20mm,bottom=22mm}
\usepackage{tikz}
\usetikzlibrary{arrows.meta,positioning,shapes.geometric}
\usepackage{xcolor}
\usepackage{tabularx,array,booktabs}
\usepackage{fancyhdr}
\usepackage{enumitem}
\usepackage{microtype}
\usepackage{needspace}
\usepackage{setspace}
\usepackage{hyperref}
\hypersetup{hidelinks,pdftitle={Каноническое разложение составного числа на простые множители},pdfauthor={Учебно-методическое пособие}}
\definecolor{ink}{HTML}{253644}
\definecolor{navy}{HTML}{324F69}
\definecolor{slate}{HTML}{798E9F}
\definecolor{muted}{HTML}{586B78}
\definecolor{pale}{HTML}{EDF2F5}
\color{ink}
\setstretch{1.08}
\setlength{\parindent}{0pt}
\setlength{\parskip}{4.2pt}
\setlist[itemize]{leftmargin=6mm,itemsep=2pt,topsep=3pt}
\setcounter{secnumdepth}{0}
\pagestyle{fancy}
\fancyhf{}
\fancyhead[L]{\footnotesize\color{muted} ЕГЭ по математике\,|\,Целые числа}
\fancyhead[R]{\footnotesize\color{muted} Каноническое разложение}
\fancyfoot[C]{\small\color{muted}\thepage}
\renewcommand{\headrulewidth}{0.35pt}
\setlength{\headheight}{14pt}
\renewcommand{\arraystretch}{1.38}
\newcommand{\myNOD}{\mathop{\text{НОД}}\nolimits}
\newcommand{\myNOK}{\mathop{\text{НОК}}\nolimits}
\newcommand{\blocktitle}[2]{\clearpage\section{Блок #1. #2}\vspace{-2mm}\noindent\color{slate}\rule{\linewidth}{.6pt}\color{ink}\vspace{2mm}}
\newcommand{\concept}[1]{\par\needspace{4\baselineskip}{\color{navy}\bfseries #1.}\enspace}
\newcommand{\example}[1]{\par\needspace{6\baselineskip}{\color{navy}\bfseries Разобранный пример.} #1\par}
\newcommand{\warning}[1]{\par\needspace{4\baselineskip}{\color{muted}\bfseries Типичная ошибка.} #1\par}
\newcommand{\selfcheck}[1]{\par\needspace{3\baselineskip}{\color{muted}\bfseries Проверка себя.} #1\par}
\newcommand{\minitrain}{\par\needspace{7\baselineskip}\vspace{1mm}{\color{navy}\bfseries Три мини-упражнения}\par\vspace{-1mm}}
\newcommand{\mini}[2]{\par\noindent\hspace*{2mm}\textbf{#1)}\enspace #2\par}
\newcommand{\task}[2]{\par\needspace{2\baselineskip}\noindent\makebox[8mm][l]{\color{navy}\bfseries #1.}#2\par\vspace{1.8mm}}
\newcommand{\training}[2]{\clearpage\section*{Тренировочный блок #1}\textcolor{muted}{#2}\par\vspace{3mm}\rule{\linewidth}{.4pt}\vspace{2mm}}
\newcommand{\answers}[2]{\clearpage\section*{Ответы к тренировочному блоку #1}\vspace{3mm}#2}
\newcommand{\anstable}[1]{\begin{center}\small\begin{tabularx}{\linewidth}{>{\centering\arraybackslash}p{9mm}X>{\centering\arraybackslash}p{9mm}X}\toprule\textbf{№}&\textbf{Ответ}&\textbf{№}&\textbf{Ответ}\\\midrule #1\\\bottomrule\end{tabularx}\end{center}}
\tikzset{
 st/.style={ellipse,draw=navy,line width=.8pt,fill=pale,align=center,minimum height=10mm,text width=43mm,font=\small},
 pr/.style={rectangle,rounded corners=2mm,draw=slate,line width=.8pt,align=center,minimum height=12mm,text width=66mm,font=\small},
 de/.style={diamond,aspect=2.4,draw=navy,line width=.8pt,align=center,text width=33mm,inner sep=1.8mm,font=\small},
 ar/.style={-{Stealth[length=2.4mm,width=1.7mm]},line width=.8pt,draw=navy},
 ann/.style={font=\footnotesize\bfseries,inner sep=1.5pt,color=navy}
}
\newcommand{\schemetitle}[1]{\clearpage\section*{Алгоритм: #1}\vspace{11mm}}
\begin{document}
\begin{titlepage}
\thispagestyle{empty}
\vspace*{24mm}
{\color{navy}\large УЧЕБНО-МЕТОДИЧЕСКОЕ ПОСОБИЕ}\par
\vspace{12mm}
{\fontsize{25}{33}\selectfont\bfseries Каноническое разложение\\[2mm] составного числа\\[2mm] на простые множители}\par
\vspace{9mm}
{\Large Теория чисел. Подготовка к профильному ЕГЭ\par}
\vspace{17mm}
{\color{muted}\large От простых множителей к делителям,\\[2mm] НОД, НОК и доказательным задачам}\par
\vspace{18mm}
\begin{tikzpicture}[x=1cm,y=1cm]
\draw[navy,line width=.95pt] (0,0) .. controls (2.9,.9) and (6.0,-.65) .. (8.3,.14) .. controls (10.7,.75) and (12.8,.28) .. (15.6,.48);
\draw[slate,line width=.35pt] (0,-.36) -- (15.6,-.36);
\end{tikzpicture}
\vfill
\end{titlepage}
\clearpage
\section*{Карта темы}
\textbf{Ключевая идея:} натуральное число больше единицы полностью характеризуется тем, какие простые числа входят в его разложение и с какими показателями.
\begin{center}
\begin{tikzpicture}[node distance=7.5mm]
\node[pr,text width=111mm] (a) {\textbf{1. Раскладываем}\quad простые делители, степени и единственность};
\node[pr,text width=111mm,below=of a] (b) {\textbf{2. Читаем разложение}\quad делимость и натуральные делители};
\node[pr,text width=111mm,below=of b] (c) {\textbf{3. Сравниваем числа}\quad НОД, НОК и взаимная простота};
\node[pr,text width=111mm,below=of c] (d) {\textbf{4. Исследуем структуру}\quad сумма делителей, квадраты и кубы};
\draw[ar](a)--(b);\draw[ar](b)--(c);\draw[ar](c)--(d);
\end{tikzpicture}
\end{center}
\concept{Условные обозначения} $\mathbb N=\{1,2,3,\ldots\}$; $d\mid n$ --- число $d$ делит $n$; $\myNOD(a,b)$ и $\myNOK(a,b)$ --- наибольший общий делитель и наименьшее общее кратное. Обозначения $\tau(n)$ и $\sigma(n)$ используются для количества и суммы \emph{положительных} делителей числа $n$.
\concept{Маршрут решения} Распознать задачу $\rightarrow$ получить каноническое разложение $\rightarrow$ выбрать показатели степеней $\rightarrow$ применить правило $\rightarrow$ проверить полноту ответа.
\concept{Границы} Для $n=1$ простых множителей нет; при этом $\tau(1)=1$ и $\sigma(1)=1$. Для отрицательного целого числа сначала раскладывают его модуль, затем восстанавливают знак «минус».

% БЛОК 1
\blocktitle{1}{Как устроено каноническое разложение}
\subsection{1.1. Простые множители и способ разложения}
\concept{Цель} Научиться записывать составное число в виде произведения степеней \emph{различных} простых чисел.
\concept{Определение} Каноническим разложением натурального $n>1$ называют представление
\[
 n=p_1^{\alpha_1}p_2^{\alpha_2}\cdots p_k^{\alpha_k},
 \qquad p_1<p_2<\cdots<p_k,
 \qquad \alpha_i\in\mathbb N,
\]
где все $p_i$ --- простые числа. Показатель $\alpha_i$ показывает, сколько раз множитель $p_i$ встречается в произведении.
\concept{Сравнение} Запись $360=2\cdot2\cdot2\cdot3\cdot3\cdot5$ верна, но каноническая запись компактнее: $360=2^3\cdot3^2\cdot5$. Запись $360=8\cdot9\cdot5$ \emph{не каноническая}, так как $8$ и $9$ составные.
\concept{Почему разложение существует} Если $n$ простое, разложение состоит из самого $n$. Если $n$ составное, $n=ab$, где $1<a,b<n$. Продолжаем раскладывать составные множители. На каждом этапе множители уменьшаются, поэтому процесс заканчивается простыми числами.
\example{Разложим $756$, начиная с наименьшего простого делителя:
\begin{align*}
756&=2\cdot378=2^2\cdot189,\\
189&=3\cdot63=3^2\cdot21=3^3\cdot7.
\end{align*}
Следовательно,
\[
 \boxed{756=2^2\cdot3^3\cdot7}.
\]
Проверка: $4\cdot27\cdot7=108\cdot7=756$.}
\warning{После нахождения делителя не следует прекращать разложение: например, $84=4\cdot21$ ещё не каноническая запись. Нужно разложить \emph{все} составные множители.}
\minitrain
\mini{1}{Разложите $84$ на простые множители. ($\boldsymbol{2^2\cdot3\cdot7}$.)}
\mini{2}{Разложите $150$. ($\boldsymbol{2\cdot3\cdot5^2}$.)}
\mini{3}{Разложите $231$. ($\boldsymbol{3\cdot7\cdot11}$.)}

\subsection{1.2. Единственность разложения и сравнение степеней}
\concept{Теорема (основная теорема арифметики)} Каждое натуральное число $n>1$ разлагается в произведение простых чисел \emph{единственным образом}, если не учитывать порядок множителей.
\concept{Ключевой факт} Если простое $p$ делит произведение $ab$, то $p\mid a$ или $p\mid b$. Действительно, если $p\nmid a$, то $\myNOD(p,a)=1$. По тождеству Безу найдутся целые $u,v$ такие, что $up+va=1$. Умножим на $b$: $upb+vab=b$; обе части суммы слева делятся на $p$, поэтому $p\mid b$.
\concept{Идея доказательства единственности} Допустим, у $n$ есть два разложения на простые множители. Первый простой множитель первого разложения делит всё второе произведение, следовательно, делит один из его простых множителей. Простые множители равны. Сократим общий множитель и повторим рассуждение. Совпадут все простые множители, включая кратности.
\example{Пусть
\[
 2^a\cdot3^b=432,\qquad a,b\in\mathbb N.
\]
Последовательным делением получаем $432=16\cdot27=2^4\cdot3^3$. По единственности канонического разложения $a=4$, $b=3$. Проверка: $2^4\cdot3^3=16\cdot27=432$.}
\concept{Практическое следствие} Нельзя «компенсировать» недостающую двойку лишней тройкой: два произведения степеней простых равны тогда и только тогда, когда показатели при \emph{каждом} простом совпадают.
\minitrain
\mini{1}{Запишите каноническое разложение $72$. ($\boldsymbol{2^3\cdot3^2}$.)}
\mini{2}{Запишите каноническое разложение $180$. ($\boldsymbol{2^2\cdot3^2\cdot5}$.)}
\mini{3}{Если $2^a3^b=108$ и $a,b\in\mathbb N$, найдите $a,b$. ($\boldsymbol{a=2,\ b=3}$.)}

\schemetitle{разложение натурального числа на простые множители}
\begin{center}
\begin{tikzpicture}[node distance=11mm]
\node[st] (s) {Дано $n>1$};
\node[pr,below=of s] (a) {Положите $m=n$; список пуст};
\node[de,below=16mm of a] (q) {$m=1$?};
\node[pr,text width=40mm,right=34mm of q] (end) {Сгруппируйте равные простые множители};
\node[pr,below=18mm of q] (b) {Найдите наименьший простой делитель $p$};
\node[pr,below=of b] (c) {Запишите множитель $p$; замените $m$ на $m/p$};
\draw[ar](s)--(a);\draw[ar](a)--(q);
\draw[ar](q.east)--node[ann,above]{да}(end.west);
\draw[ar](q.south)--node[ann,right]{нет}(b.north);
\draw[ar](b)--(c);
\draw[ar](c.west)--++(-23mm,0)|-(q.west);
\end{tikzpicture}
\end{center}

\training{1}{Разложение, кратности простых множителей и единственность записи.}
\task{1}{Найдите каноническое разложение числа $90$.}
\task{2}{Найдите каноническое разложение числа $196$.}
\task{3}{Разложите $420$ на простые множители.}
\task{4}{Разложите $945$ на простые множители.}
\task{5}{Разложите $756$ на простые множители.}
\task{6}{Вычислите натуральное число, записанное в виде $2^3\cdot3^2\cdot5$.}
\task{7}{Равны ли числа $2^4\cdot3^2\cdot5$ и $2^3\cdot3^3\cdot5$? Ответ обоснуйте через показатели.}
\task{8}{Найдите натуральные показатели $a,b$ из равенства $2^a3^b=432$.}
\task{9}{Известно, что $864=2^a3^b$, где $a,b\in\mathbb N$. Найдите $a$ и $b$.}
\task{10}{Найдите наименьшее натуральное число $N$, кратное $12$ и имеющее ровно три \emph{различных} простых делителя. Объясните минимальность.}
\answers{1}{\anstable{
1&$2\cdot3^2\cdot5$&6&$360$\\
2&$2^2\cdot7^2$&7&Нет: $720\ne1080$\\
3&$2^2\cdot3\cdot5\cdot7$&8&$a=4$, $b=3$\\
4&$3^3\cdot5\cdot7$&9&$a=5$, $b=3$\\
5&$2^2\cdot3^3\cdot7$&10&$60=2^2\cdot3\cdot5$
}}
\selfcheck{Для задания 10: делимость на $12=2^2\cdot3$ требует простые $2$ и $3$; наименьший третий простой множитель --- $5$, откуда $N\ge12\cdot5=60$.}

% БЛОК 2
\blocktitle{2}{Делимость и количество натуральных делителей}
\subsection{2.1. Делимость как сравнение показателей}
\concept{Цель} Решать вопросы о делимости без вычисления больших чисел.
\concept{Критерий делимости} Пусть все простые множители чисел $a$ и $b$ включены в единый список $p_1,\ldots,p_k$, а отсутствующим простым соответствуют нулевые показатели:
\[
 a=\prod_{i=1}^k p_i^{\alpha_i},\qquad
 b=\prod_{i=1}^k p_i^{\beta_i}.
\]
Тогда
\[
 \boxed{a\mid b\ \Longleftrightarrow\ \alpha_i\le\beta_i\text{ для каждого }i.}
\]
\concept{Почему так} Если $a\mid b$, то $b=ac$ и показатели в $b$ получаются сложением показателей в $a$ и $c$. Обратно, если все $\beta_i-\alpha_i\ge0$, то $b/a=\prod p_i^{\beta_i-\alpha_i}$ --- натуральное число.
\example{Проверим, делится ли $2160$ на $72$:
\[
2160=2^4\cdot3^3\cdot5,\qquad72=2^3\cdot3^2.
\]
Показатели $4\ge3$ и $3\ge2$, значит, делится. Частное равно
\[
 \frac{2160}{72}=2^{4-3}\cdot3^{3-2}\cdot5=30.
\]
Напротив, $2160\nmid72$: нужно различать направление делимости.}
\warning{При проверке $a\mid b$ сравнивают показатели у \emph{делителя} $a$ с показателями у \emph{делимого} $b$; направление неравенств нельзя менять.}
\minitrain
\mini{1}{Верно ли $2^3\cdot3^2\mid2^4\cdot3^2\cdot7$? (\textbf{Да}.)}
\mini{2}{Какие из $12,16,18$ делят $360$? ($\boldsymbol{12,18}$; число $16$ не делит $360$.)}
\mini{3}{При каких целых $a\ge0$ число $2^a3^2$ делит $72$? ($\boldsymbol{a=0,1,2,3}$.)}

\subsection{2.2. Как перечислять и считать делители}
\concept{Теорема} Если $n=p_1^{\alpha_1}\cdots p_k^{\alpha_k}$, то всякий его положительный делитель имеет вид
\[
 d=p_1^{\beta_1}\cdots p_k^{\beta_k},\qquad
 0\le\beta_i\le\alpha_i.
\]
Число таких делителей равно
\[
 \boxed{\tau(n)=(\alpha_1+1)\cdots(\alpha_k+1).}
\]
\concept{Объяснение формулы} Для $p_1$ есть $\alpha_1+1$ вариантов показателя: $0,1,\ldots,\alpha_1$. Для $p_2$ --- $\alpha_2+1$ вариантов и так далее. Выборы независимы, поэтому количества вариантов перемножаются.
\example{Пусть $n=72=2^3\cdot3^2$. Каждый делитель задаётся парой $(\beta,\gamma)$, где $\beta\in\{0,1,2,3\}$ и $\gamma\in\{0,1,2\}$. Поэтому $\tau(72)=4\cdot3=12$.
\begin{center}
\small\begin{tabular}{c|ccc}
\toprule
&$3^0$&$3^1$&$3^2$\\\midrule
$2^0$&$1$&$3$&$9$\\
$2^1$&$2$&$6$&$18$\\
$2^2$&$4$&$12$&$36$\\
$2^3$&$8$&$24$&$72$\\\bottomrule
\end{tabular}
\end{center}
Таблица перечисляет все $12$ делителей ровно по одному разу.}
\concept{Делители с дополнительным условием} Для делителей, кратных $12=2^2\cdot3$, показатели двойки и тройки должны быть не ниже $2$ и $1$. Например, у $n=2^4\cdot3^2\cdot5$ таких делителей $3\cdot2\cdot2=12$: для двойки доступны степени $2,3,4$; для тройки --- $1,2$; для пятёрки --- $0,1$.
\concept{Полезный частный случай} Если $n=p^a$ для простого $p$, то $\tau(n)=a+1$. Если $n=pq$ для различных простых $p,q$, то ровно четыре делителя: $1,p,q,pq$.
\selfcheck{Не забывайте делитель $1$: ему соответствует выбор \emph{всех} показателей, равных нулю. Само число $n$ тоже является делителем.}
\minitrain
\mini{1}{Найдите $\tau(2^3\cdot5^2)$. ($\boldsymbol{12}$.)}
\mini{2}{Сколько положительных делителей у $45$? ($\boldsymbol{6}$.)}
\mini{3}{Найдите $\tau(2^3\cdot3)$. ($\boldsymbol{8}$.)}

\schemetitle{подсчёт натуральных делителей}
\begin{center}
\begin{tikzpicture}[node distance=11mm]
\node[st] (s) {Дано натуральное $n$};
\node[de,below=16mm of s] (q) {$n=1$?};
\node[pr,text width=30mm,right=41mm of q] (o) {$\tau(1)=1$};
\node[pr,below=19mm of q] (a) {Разложите $n$ на простые множители};
\node[pr,below=of a] (b) {Выпишите показатели $\alpha_1,\ldots,\alpha_k$};
\node[pr,below=of b] (c) {Вычислите $\prod(\alpha_i+1)$};
\node[st,below=of c] (e) {Получено $\tau(n)$};
\draw[ar](s)--(q);
\draw[ar](q.east)--node[ann,above]{да}(o.west);
\draw[ar](q.south)--node[ann,right]{нет}(a.north);
\draw[ar](a)--(b);\draw[ar](b)--(c);\draw[ar](c)--(e);
\end{tikzpicture}
\end{center}

\training{2}{Делимость, выбор показателей и комбинаторика делителей.}
\task{1}{Сколько натуральных делителей у $48$?}
\task{2}{Сколько натуральных делителей у $75$?}
\task{3}{Перечислите все положительные делители числа $30$.}
\task{4}{Какое из чисел $2^5\cdot3^2$ и $2^4\cdot3^3$ делит другое?}
\task{5}{Найдите число \emph{чётных} натуральных делителей числа $2^4\cdot3^2$.}
\task{6}{Сколько натуральных делителей числа $2^4\cdot3^3\cdot5$ кратны $12$?}
\task{7}{Найдите все пары натуральных показателей $(a,b)$, для которых $\tau(2^a3^b)=12$.}
\task{8}{Найдите все натуральные числа от $1$ до $100$ включительно, имеющие ровно девять положительных делителей.}
\task{9}{При каком целом $a\ge0$ число $2^a\cdot3^2\cdot7$ имеет ровно $30$ натуральных делителей?}
\task{10}{Сколько положительных делителей числа $2^5\cdot3^3\cdot5^2$ являются полными квадратами?}
\answers{2}{\anstable{
1&$10$&6&$18$\\
2&$6$&7&$(1,5),(2,3),(3,2),(5,1)$\\
3&$1,2,3,5,6,10,15,30$&8&$36,100$\\
4&Ни одно не делит другое&9&$a=4$\\
5&$12$&10&$12$
}}
\selfcheck{Для задания 8: девять делителей возможны только для $p^8$ либо $p^2q^2$ при $p\ne q$. На отрезке $[1,100]$ числа первого вида отсутствуют, второго --- $36$ и $100$.}

% БЛОК 3
\blocktitle{3}{НОД, НОК и взаимная простота}
\subsection{3.1. Минимальные и максимальные показатели}
\concept{Цель} По каноническим разложениям двух чисел находить их наибольший общий делитель и наименьшее общее кратное.
\concept{Теорема} Пусть $a=\prod p_i^{\alpha_i}$, $b=\prod p_i^{\beta_i}$, причём отсутствующие множители учтены показателем $0$. Тогда
\begin{align*}
\myNOD(a,b)&=\prod p_i^{\min(\alpha_i,\beta_i)},\\
\myNOK(a,b)&=\prod p_i^{\max(\alpha_i,\beta_i)}.
\end{align*}
\concept{Смысл} Общий делитель не может содержать простого множителя больше, чем каждое из чисел, поэтому берём \emph{минимальные} степени. Общее кратное должно содержать степени, достаточные для обоих чисел, поэтому берём \emph{максимальные}.
\example{Найдём НОД и НОК чисел $144$ и $210$:
\begin{align*}
144&=2^4\cdot3^2,\\
210&=2\cdot3\cdot5\cdot7.
\end{align*}
Тогда
\begin{align*}
\myNOD(144,210)&=2\cdot3=6,\\
\myNOK(144,210)&=2^4\cdot3^2\cdot5\cdot7=5040.
\end{align*}
\textbf{Проверка:} $6\cdot5040=30240=144\cdot210$.}
\concept{Полезное тождество} Для всех натуральных $a,b$ верно
\[
 \boxed{\myNOD(a,b)\cdot\myNOK(a,b)=ab.}
\]
Оно следует из равенства $\min(\alpha,\beta)+\max(\alpha,\beta)=\alpha+\beta$ для каждого простого.
\minitrain
\mini{1}{Найдите $\myNOD(72,120)$ и $\myNOK(72,120)$. ($\boldsymbol{24,360}$.)}
\mini{2}{Найдите $\myNOD(2^4\cdot3^2,2^2\cdot3^3)$. ($\boldsymbol{36}$.)}
\mini{3}{Найдите $\myNOK(18,30)$. ($\boldsymbol{90}$.)}

\subsection{3.2. Взаимно простые числа и восстановление пар}
\concept{Определение} Два натуральных числа взаимно просты, если их НОД равен $1$. В канонических разложениях таких чисел \emph{нет общих простых множителей}.
\concept{Следствия} Для взаимно простых $a,b$ выполнено $\myNOK(a,b)=ab$. Если $a\mid bc$ и $\myNOD(a,b)=1$, то $a\mid c$: все простые множители числа $a$ вынужденно находятся в $c$ с нужными кратностями.
\example{Пусть $\myNOD(a,b)=12$ и $\myNOK(a,b)=180$. Найдём все \emph{неупорядоченные} пары $(a,b)$.
Запишем $a=12u$, $b=12v$, где $\myNOD(u,v)=1$. По формуле НОД--НОК получаем
\[
12u\cdot12v=12\cdot180\quad\Longrightarrow\quad uv=15.
\]
Взаимно простые пары натуральных множителей числа $15$ (без учёта порядка): $(1,15)$ и $(3,5)$. Следовательно,
\[
 \boxed{\{a,b\}=\{12,180\}\quad\text{или}\quad\{a,b\}=\{36,60\}.}
\]
Проверка: у обеих пар НОД равен $12$, а НОК --- $180$.}
\warning{Из $ab=\myNOD(a,b)\cdot\myNOK(a,b)$ нельзя заключить, что пара $(a,b)$ определена однозначно.}
\minitrain
\mini{1}{Взаимно ли просты $35$ и $64$? (\textbf{Да}.)}
\mini{2}{Найдите $\myNOD(2^3\cdot3\cdot5,2\cdot7^2)$. ($\boldsymbol{2}$.)}
\mini{3}{Приведите взаимно простые $a,b>1$ с произведением $180$. ($\boldsymbol{4\text{ и }45}$.)}

\schemetitle{нахождение НОД и НОК по простым множителям}
\begin{center}
\begin{tikzpicture}[node distance=11mm]
\node[st] (s) {Даны натуральные $a,b$};
\node[pr,below=of s] (a) {Разложите $a$ и $b$ на простые множители};
\node[pr,below=of a] (b) {Объедините списки простых; допишите нулевые показатели};
\node[de,below=17mm of b] (q) {Найти НОД или НОК?};
\node[pr,text width=42mm,below=23mm of q,xshift=-46mm] (l) {НОД: возьмите \emph{минимальные} показатели};
\node[pr,text width=42mm,below=23mm of q,xshift=46mm] (r) {НОК: возьмите \emph{максимальные} показатели};
\node[st,text width=45mm,below=58mm of q] (e) {Перемножьте степени};
\draw[ar](s)--(a);\draw[ar](a)--(b);\draw[ar](b)--(q);
\draw[ar](q.west)-|(l.north);
\draw[ar](q.east)-|(r.north);
\node[ann] at ([xshift=-12mm,yshift=2mm]q.west) {НОД};
\node[ann] at ([xshift=12mm,yshift=2mm]q.east) {НОК};
\draw[ar](l.south)|-(e.west);\draw[ar](r.south)|-(e.east);
\end{tikzpicture}
\end{center}

\training{3}{НОД, НОК, показатели, взаимная простота и восстановление чисел.}
\task{1}{Найдите $\myNOD(84,126)$.}
\task{2}{Найдите $\myNOK(24,90)$.}
\task{3}{Найдите $\myNOD(2^4\cdot3^2,2^2\cdot3^5\cdot7)$.}
\task{4}{Найдите $\myNOK(2^2\cdot5^3,2^4\cdot3\cdot5)$.}
\task{5}{Найдите НОД и НОК чисел $180$ и $252$.}
\task{6}{Какие пары взаимно просты: $(35,64)$, $(45,60)$, $(49,81)$?}
\task{7}{Пусть $a\ge0$ --- целое и $\myNOD(2^a\cdot3,2^2\cdot3^3)=6$. Найдите $a$.}
\task{8}{Найдите все неупорядоченные пары $(a,b)$ с $\myNOD(a,b)=12$ и $\myNOK(a,b)=180$.}
\task{9}{Докажите, что из $\myNOD(a,b)=1$ следует $\myNOD(a,a+b)=1$.}
\task{10}{Найдите наименьшее натуральное $n$, для которого $\myNOK(n,18)=180$.}
\answers{3}{\anstable{
1&$42$&6&$(35,64)$ и $(49,81)$\\
2&$360$&7&$a=1$\\
3&$36$&8&$\{12,180\}$ и $\{36,60\}$\\
4&$6000$&9&Общий делитель делит $b=(a+b)-a$\\
5&НОД $36$, НОК $1260$&10&$n=20$
}}
\selfcheck{Для задания 10: $180=2^2\cdot3^2\cdot5$, $18=2\cdot3^2$. Число $n$ обязано содержать $2^2$ и $5$, а наименьшее такое $n=20$; $\myNOK(20,18)=180$.}

% БЛОК 4
\blocktitle{4}{Сумма делителей и полные степени}
\subsection{4.1. Сумма всех натуральных делителей}
\concept{Цель} Находить сумму делителей, не перебирая их вручную.
\concept{Теорема} Если $n=p_1^{\alpha_1}\cdots p_k^{\alpha_k}$, то сумма всех \emph{положительных} делителей равна
\[
\boxed{\sigma(n)=\prod_{i=1}^k(1+p_i+p_i^2+\cdots+p_i^{\alpha_i}).}
\]
Для вычислений можно использовать формулу геометрической прогрессии:
\[
1+p+\cdots+p^\alpha=\frac{p^{\alpha+1}-1}{p-1}.
\]
\concept{Почему формула работает} Раскрытие скобок в произведении выбирает из каждой скобки ровно одну степень $p_i^{\beta_i}$, $0\le\beta_i\le\alpha_i$. Каждому набору показателей соответствует ровно один натуральный делитель числа $n$, поэтому сумма включает каждый делитель по одному разу.
\example{Найдём сумму делителей $72=2^3\cdot3^2$:
\begin{align*}
\sigma(72)&=(1+2+2^2+2^3)(1+3+3^2)\\
&=(1+2+4+8)(1+3+9)=15\cdot13=195.
\end{align*}
\textbf{Проверка:} из таблицы делителей в блоке 2 имеем $1+2+3+4+6+8+9+12+18+24+36+72=195$.}
\concept{Сравнение} Формула $\tau(n)$ \emph{считает} варианты показателей, а формула $\sigma(n)$ \emph{суммирует} соответствующие степени. Эти формулы нельзя подменять друг другом.
\minitrain
\mini{1}{Найдите $\sigma(36)$. ($\boldsymbol{91}$.)}
\mini{2}{Найдите $\sigma(2^3)$. ($\boldsymbol{15}$.)}
\mini{3}{Найдите $\sigma(15)$. ($\boldsymbol{24}$.)}

\subsection{4.2. Когда число является квадратом или кубом}
\concept{Критерий полной степени} Натуральное число $n=\prod p_i^{\alpha_i}$ является точной $r$-й степенью натурального числа ($r\ge2$) тогда и только тогда, когда \emph{каждый} показатель $\alpha_i$ делится на $r$.
\concept{Почему} Если $n=m^r$, то показатели простых множителей числа $n$ в $r$ раз больше показателей в $m$. И наоборот, если все $\alpha_i$ кратны $r$, то $n=(\prod p_i^{\alpha_i/r})^r$.
\example{Найдём наименьшее натуральное $k$, для которого $180k$ --- квадрат. Разложим
\[
180=2^2\cdot3^2\cdot5.
\]
Показатель при $5$ нечётный, а при $2,3$ --- чётные. Достаточно и необходимо умножить на $5$. Получим
\[
180\cdot5=900=30^2.
\]
Множитель $k=5$ минимален: любой подходящий множитель обязан содержать хотя бы одну пятёрку.}
\concept{Как узнать квадрат по числу делителей} Если $n$ --- полный квадрат, то $\tau(n)$ нечётно: все показатели $\alpha_i$ чётны, значит, все $\alpha_i+1$ нечётны. Верно и обратное: если произведение $\prod(\alpha_i+1)$ нечётно, то все показатели чётны. Следовательно,
\[
\boxed{\tau(n)\text{ нечётно}\ \Longleftrightarrow\ n\text{ --- полный квадрат}.}
\]
\concept{Сколько делителей являются квадратами} У квадратного делителя все показатели чётные. Для каждого $p_i$ можно выбрать $0,2,4,\ldots,2\lfloor\alpha_i/2\rfloor$, откуда число квадратных делителей равно $\prod_i(\lfloor\alpha_i/2\rfloor+1)$. Например, для $360=2^3\cdot3^2\cdot5$ их $(1+1)(1+1)(0+1)=4$: $1,4,9,36$.
\warning{Нельзя проверять только один показатель. Например, $2^2\cdot3^3$ \emph{не является квадратом}, хотя степень двойки чётная.}
\minitrain
\mini{1}{Является ли $2^4\cdot3^2$ квадратом? (\textbf{Да}: $12^2$.)}
\mini{2}{Является ли $2^3\cdot3^2$ кубом? (\textbf{Нет}.)}
\mini{3}{Найдите наименьшее $k\in\mathbb N$, при котором $72k$ --- квадрат. ($\boldsymbol{k=2}$.)}

\schemetitle{дополнение разложения до полного квадрата или куба}
\begin{center}
\begin{tikzpicture}[node distance=11mm]
\node[st] (s) {Даны $n>1$ и степень $r\ge2$};
\node[pr,below=of s] (a) {Разложите $n=\prod p_i^{\alpha_i}$};
\node[de,below=17mm of a] (q) {Все $\alpha_i$ кратны $r$?};
\node[pr,text width=34mm,right=41mm of q] (yes) {$n$ уже является $r$-й степенью};
\node[pr,below=20mm of q] (b) {Добавьте к каждому $\alpha_i$ минимум $\delta_i\ge0$ до кратности $r$};
\node[pr,below=of b] (c) {Возьмите $k=\prod p_i^{\delta_i}$};
\node[st,below=of c] (e) {$nk$ --- $r$-я степень; $k$ минимально};
\draw[ar](s)--(a);\draw[ar](a)--(q);
\draw[ar](q.east)--node[ann,above]{да}(yes.west);
\draw[ar](q.south)--node[ann,right]{нет}(b.north);
\draw[ar](b)--(c);\draw[ar](c)--(e);
\end{tikzpicture}
\end{center}

\training{4}{Сумма делителей, число делителей квадратов и дополнение до полных степеней.}
\task{1}{Найдите сумму натуральных делителей числа $12$.}
\task{2}{Найдите количество и сумму положительных делителей числа $16$.}
\task{3}{Найдите $\sigma(45)$.}
\task{4}{Является ли $450$ полным квадратом? Обоснуйте по показателям.}
\task{5}{Найдите наименьшее натуральное $k$, при котором $180k$ --- полный квадрат.}
\task{6}{Найдите наименьшее натуральное $k$, при котором $72k$ --- полный куб.}
\task{7}{Сколько положительных делителей числа $360$ являются полными квадратами?}
\task{8}{Найдите натуральное $a$, если $\sigma(2^a)=63$.}
\task{9}{Докажите, что число с чётным количеством положительных делителей не может быть полным квадратом.}
\task{10}{Найдите наименьшее натуральное $k$, при котором $90k$ --- полный куб. Укажите сам полученный куб.}
\answers{4}{\anstable{
1&$28$&6&$k=3$, $72k=216=6^3$\\
2&$\tau(16)=5$, $\sigma(16)=31$&7&$4$\\
3&$78$&8&$a=5$\\
4&Нет: $450=2\cdot3^2\cdot5^2$&9&У квадрата все показатели чётные, $\tau$ нечётно\\
5&$k=5$, $180k=900=30^2$&10&$k=300$, $90k=27000=30^3$
}}
\selfcheck{Для задания 10: $90=2\cdot3^2\cdot5$. Чтобы показатели стали кратными $3$, минимально добавить $2^2\cdot3\cdot5^2=300$.}

% ИТОГОВАЯ
\clearpage
\section*{Итоговая обзорная контрольная работа}
\textcolor{muted}{10 заданий: 1--3 --- базовые; 4--7 --- среднего уровня; 8--9 --- повышенные; 10 --- комплексное.}
\par\vspace{3mm}\rule{\linewidth}{.4pt}\vspace{2mm}
\task{1}{Найдите каноническое разложение числа $540$.}
\task{2}{Сколько натуральных делителей у числа $72$?}
\task{3}{Определите по каноническому разложению, является ли число $441$ квадратом и является ли оно кубом.}
\task{4}{Верно ли, что $2^3\cdot3^2\cdot5$ делит $2^4\cdot3^3\cdot5^2$? Обоснуйте.}
\task{5}{Найдите НОД и НОК чисел $420$ и $630$.}
\task{6}{Сколько натуральных делителей числа $2^5\cdot3^3$ кратны $24$?}
\task{7}{Найдите сумму всех натуральных делителей числа $84$.}
\task{8}{Известно, что число $n=2^a\cdot3^2$, $a\in\mathbb N$, имеет ровно $18$ натуральных делителей. Найдите $n$.}
\task{9}{Докажите, что не существует натурального полного квадрата с ровно $14$ положительными делителями.}
\task{10}{Найдите \emph{наименьшее} натуральное $n$, кратное $12$ и имеющее ровно $24$ положительных делителя. Докажите минимальность найденного числа.}
\clearpage
\section*{Ответы к итоговой обзорной контрольной работе}
\anstable{
1&$540=2^2\cdot3^3\cdot5$&6&$9$\\
2&$12$&7&$224$\\
3&Квадрат $21^2$; не куб&8&$288$ ($a=5$)\\
4&Да: все показатели не больше&9&У квадрата $\tau(n)$ нечётно\\
5&НОД $210$, НОК $1260$&10&$360=2^3\cdot3^2\cdot5$
}
\vspace{4mm}
\concept{Пояснение к задаче 10} Число $360$ кратно $12$ и имеет $(3+1)(2+1)(1+1)=24$ делителя. Для минимальности рассмотрим все формы разложения с $24$ делителями.
\begin{itemize}
\item При двух различных простых показатели дают $24=(a+1)(b+1)$ с $a\ge2$ для кратности $12$. Допустимые пары показателей $(a,b)$: $(2,7)$, $(3,5)$, $(5,3)$, $(7,2)$, $(11,1)$. Наименьшее соответствующее число равно $2^5\cdot3^3=864>360$.
\item При трёх различных простых минимальная допустимая комбинация показателей --- $(3,2,1)$, дающая $2^3\cdot3^2\cdot5=360$. Другая форма $(5,1,1)$ даёт не менее $2^5\cdot3\cdot5=480$.
\item При четырёх различных простых $24=3\cdot2\cdot2\cdot2$, поэтому минимально возможное число, кратное $12$, не меньше $2^2\cdot3\cdot5\cdot7=420$.
\item С одним простым число, кратное $12$, получить нельзя; с пятью и более различными простыми множителями количество делителей не менее $2^5=32$.
\end{itemize}
Значит, наименьшее подходящее число действительно равно $360$.

\clearpage
\section*{Памятка по каноническому разложению}
\begin{center}
\begin{tabularx}{\linewidth}{>{\raggedright\arraybackslash}p{50mm}X}
\toprule
\textbf{Вопрос}&\textbf{Что использовать}\\\midrule
Как записать число?&$n=\prod p_i^{\alpha_i}$, $p_i$ --- различные простые\\
Когда $a\mid b$?&$\alpha_i\le\beta_i$ для всех простых $p_i$\\
Сколько делителей?&$\tau(n)=\prod(\alpha_i+1)$\\
Как найти НОД?&$\prod p_i^{\min(\alpha_i,\beta_i)}$\\
Как найти НОК?&$\prod p_i^{\max(\alpha_i,\beta_i)}$\\
Как найти сумму делителей?&$\sigma(n)=\prod(1+p_i+\cdots+p_i^{\alpha_i})$\\
Когда число --- квадрат?&Все показатели $\alpha_i$ чётные\\
Когда число --- куб?&Все показатели $\alpha_i$ кратны $3$\\\bottomrule
\end{tabularx}
\end{center}
\concept{Самопроверка} Разложение должно содержать только простые множители. При подсчёте делителей учтите $1$ и само число. При поиске всех вариантов проверьте, что перечислены \emph{все} возможные наборы показателей. В доказательной задаче предъявленного примера недостаточно для обоснования минимальности или единственности.
\end{document}